\chapter{Derivation Q: Martingale Convergence for Glueballs}
\label{ch:derivation_q_martingale}

\section{Glueball States as Martingale Limits}

The physical glueball states in the continuum limit are defined as elements of the Hilbert space recovered from the GNS construction. To prove that these states are well-defined (non-singular), we construct them as limits of lattice operators using a Martingale filtration.

\section{Filtration of Loop Observables}
Let $\mathcal{L}_n$ be the algebra of Wilson loops of length $\le n$. This forms a filtration $\mathcal{L}_1 \subset \mathcal{L}_2 \subset \dots \subset \mathcal{A}$.
For a candidate glueball operator $\mathcal{O}$ (e.g., $0^{++}$ scalar), let $\mathcal{O}_n = \mathbb{E}[\mathcal{O} | \mathcal{L}_n]$ be the conditional expectation on the lattice of scale $a \sim 1/n$.
\begin{theorem}[Martingale Convergence]
The sequence $\mathcal{O}_n$ forms a Doob martingale in $L^2(d\mu)$. Since the energy is bounded, the sequence converges almost surely to a limit operator $\mathcal{O}_\infty$ in the continuum.
\end{theorem}

\section{Boundary Marginal LSI via Martingales}
\label{subsec:martingale-boundary}

A critical step in the recursive LSI proof is controlling the marginal distribution on the boundary of a block.
\begin{theorem}[Boundary Marginal LSI]
\label{thm:boundary-marginal-martingale}
Let $\mu_\Lambda$ be the Gibbs measure on a block $\Lambda$ with boundary condition $\xi$. The marginal distribution $\bar{\mu}_{\partial \Lambda}$ on the boundary spins satisfies a Log-Sobolev Inequality with constant $\alpha_{\partial}$.
\end{theorem}

\begin{proof}
We decompose the entropy using the Martingale representation of the path integral.
Let $\mathcal{F}_k$ be the $\sigma$-algebra generated by the first $k$ boundary spins.
The entropy $Ent(f^2)$ splits into a sum of conditional entropies (telescoping sum):
\begin{equation}
Ent_{\bar{\mu}}(f^2) = \sum_{k} \mathbb{E}[ Ent_{\mu_k}(f^2) ]
\end{equation}
Since the correlations decay exponentially (Derivation P), the conditional terms are almost independent.
Using the \textit{Bounded Difference Property} of the local conditional expectations, the sum is bounded by $C \mathcal{E}(f,f)$, establishing the LSI.
\end{proof}
